\documentclass{report}
%----- PACKAGES------------- 
\usepackage[T1]{fontenc}
\usepackage[utf8]{inputenc}
\usepackage{amsmath,amssymb,amsthm}
\usepackage[shortlabels]{enumitem}

\usepackage{graphics}
\usepackage{graphicx}
\usepackage{tikz-cd}
\usepackage{tikz}
\usetikzlibrary{backgrounds,decorations.pathreplacing,patterns,positioning,arrows.meta}
\usepackage{hyperref}
\usepackage{multicol}
\setlength{\columnsep}{5cm}
\usepackage{soul}
\usepackage{xcolor}
%----commands -------------

\newcommand{\NN}{\mathbb{N}}
\newcommand{\ZZ}{\mathbb{Z}}
\newcommand{\QQ}{\mathbb{Q}}
\newcommand{\RR}{\mathbb{R}}
\newcommand{\CC}{\mathbb{C}}



%----Theoremstyles------


\theoremstyle{plain}
\newtheorem{theorem}{Théorème}
\newtheorem{lemme}[theorem]{Lemme}
\newtheorem*{proposition}{Proposition}
\newtheorem{corollary}[theorem]{Corollary}


\theoremstyle{remark}
\newtheorem*{remarque}{Remarque}

\theoremstyle{definition}
\newtheorem*{definition}{Définition}
\newtheorem*{definitions}{Définitions}
\newtheorem*{observation}{Observation}
\newtheorem{notation}{Notation}
\newtheorem*{exemple}{Exemple}
\newtheorem{exercice}{Exercice}
\newtheorem*{Exercicesup}{Exercice supplémentaire}


%--- Structure---

\usepackage{titlesec}



\newcommand{\exo}[1]{\noindent \textbf{Exercice #1}}
\newcommand{\solution}{\noindent \textbf{Solution}}
\newcommand{\ajd}[1]{\noindent \textbf{Au programme:} #1}





\title{Box Problem}
\author{Dries Rooryck}

\begin{document}
\maketitle

6. Dans une version simplifiée du jeu télévisé, supposons avoir cent boites contenant une somme de 1, 2, 3 etc jusque 100 euros. Les boites sont mélangées. A chaque manche, vous ouvrez une boite. A chaque instant, vous avez deux options : (i) vous encaissez largent repris dans la caisse et le jeu sarrête ; (ii) vous jetez la boite que vous avez en main et vous en choisissez une nouvelle parmi celles qui restent.

In a simplified version of the TV gameshow, imagine having one hundred boxes. Each box contains between 1 and 100 euros. The boxes are mixed up and randomly distributed, meaning you do not know what the distribution of money is.
Every round, you open a box. When you open a box, you have two options:

(i) You take the money that you find in this box, and the round ends. \\
(ii) You throw away the box and choose a new box from the remaining ones. \\

If you choose the box with the largest sum of money, you win. However, you lose if there was another box with more money.\\

What could be a good strategy to win the most rounds?

\section{uniform distribution}

this is a special case\\

If we assume all numbers from 1-100 have the same probabilty of appearing in a box, we can find a good strategy only for this uniform distribution

\end{document}